% Ad Quintum Fratrem 3.8 (ad-quintum-fratrem-03-08)
% Source: https://raw.githubusercontent.com/PerseusDL/canonical-latinLit/master/data/phi0474/phi058/phi0474.phi058.perseus-lat1.xml
% Fetched by scripts/fetch_latin.py

% Dateline (Perseus): Scr. Romae ex. m. Nov. a. 700 (54) .

\ciceroLetterOpener{MARCVS QUINTO FRATRI SALUTEM}

\ciceroSection{1} superiori epistulae quod respondeam nihil est, quae plena stomachi et querelarum est, quo in genere alteram quoque te scribis pridie Labieno dedisse, qui adhuc non venerat. delevit enim mihi omnem molestiam recentior epistula. tantum te et moneo et rogo ut in istis molestiis et laboribus et desideriis recordere consilium nostrum quod fuerit profectionis tuae. non enim commoda quaedam sequebamur parva ac mediocria. quid enim erat quod discessu nostro emendum putaremus? praesidium firmissimum petebamus ex optimi et potentissimi viri benevolentia ad omnem statum nostrae dignitatis. plura ponuntur in spe quam petimus ; reliqua ad iacturam reserventur. qua re, si crebro referes animum tuum ad rationem et veteris consili nostri et spei, facilius istos militiae labores ceteraque quae te offendunt feres et tamen cum voles depones. sed eius rei maturitas nequedum venit et tamen iam adpropinquat.

\ciceroSection{2} etiam illud te admoneo ne quid ullis litteris committas, quod si prolatum sit, moleste feramus. multa sunt quae ego nescire malo quam cum aliquo periculo fieri certior. plura ad te vacuo animo scribam cum, ut spero, se Cicero meus belle habebit. tu velim cures ut sciam quibus nos dare oporteat eas quas ad te deinde litteras mittemus, Caesarisne tabellariis ut is ad te protinus mittat, an Labieni. Vbi enim isti sint Nervii et quam longe absint nescio.

\ciceroSection{3} de virtute et gravitate Caesaris, quam in summo dolore adhibuisset, magnam ex epistula tua accepi voluptatem. quod me institutum ad illum poema iubes perficere, etsi distentus cum opera tum animo sum multo magis, tamen quoniam ex epistula quam ad te miseram cognovit Caesar me aliquid esse exorsum, revertar ad institutum idque perficiam his supplicationum otiosis diebus, quibus Messalam iam nostrum reliquosque molestia levatos vehementer gaudeo, eumque quod certum consulem cum Domitio numeratis nihil a nostra opinione dissentitis. ego Messalam Caesari praestabo. sed Memmius in adventu Caesaris habet spem ; in quo illum puto errare. hic quidem friget ; Scaurum autem iam pridem Pompeius abiecit.

\ciceroSection{4} res prolatae ; ad interregnum comitia adducta. rumor dictatoris iniucundus bonis, mihi etiam magis quae loquuntur. sed tota res et timetur et refrigescit. Pompeius plane se negat velle ; antea mihi ipse non negabat Hirrus auctor fore videtur, o di, quam ineptus, quam se ipse amans sine rivali! Crassum Iunianum, hominem mihi deditum, per me deterruit velit nolit scire difficile est ; Hirro tamen agente nolle se non probabit. aliud hoc tempore de re publica nihil loquebantur ; agebatur quidem certe nihil.

\ciceroSection{5} Serrani domestici fili funus perluctuosum fuit a. d. viii K. Decembr. laudavit pater scripto meo.

\ciceroSection{6} nunc de Milone Pompeius ei nihil tribuit et omnia Cottae dicitque se perfecturum ut in illum Caesar incumbat. hoc horret Milo, nec iniuria et, si ille dictator factus sit, paene diffidit. intercessorem dictaturae si iuverit manu et praesidio suo, Pompeium metuit inimicum ; si non iuverit, timet ne per vim perferatur. ludos adparat magnificentissimos, sic, inquam, ut nemo sumptuosiores, stulte bis terque non postulatos vel quia munus magnificum dederat vel quia facultates non erant vel quia magister vel quia potuerat magistrum se non aedilem putare. omnia fere scripsi. cura , mi carissime frater, ut valeas.
